% Walkthrough recording runbook — a companion to docs/video-shotlist.md, not part of the graded report.
% Build: latexmk -pdf recording-runbook.tex   (Overleaf: Menu → Main document → docs/report/recording-runbook.tex, then Recompile)
% Standalone: no bibliography, no figures. Commands are verbatim from the README and the rehearsed shot list.

\documentclass[11pt,letterpaper]{article}
\usepackage[margin=0.75in]{geometry}
\usepackage[T1]{fontenc}
\usepackage[utf8]{inputenc}
\usepackage[expansion=false]{microtype}
\usepackage{booktabs}
\usepackage{xcolor}
\usepackage{hyperref}
\hypersetup{colorlinks=true,linkcolor=black,urlcolor=blue!50!black}
\usepackage{enumitem}
\setlist{nosep}
\usepackage{framed}
\definecolor{saybg}{RGB}{236,244,255}
\colorlet{shadecolor}{saybg}
\newcommand{\code}[1]{\texttt{#1}}
\newcommand{\say}[1]{\begin{shaded*}\noindent\textbf{Say:} #1\end{shaded*}}
\newcommand{\expect}[1]{\par\noindent\textbf{Expected:} #1\par}
\newcommand{\verify}[1]{\par\noindent\textbf{Verify before moving on:} #1\par}
\setlength{\parskip}{4pt}
\emergencystretch=3em
% Verbatim blocks in footnotesize so the README's long lines fit the text width unchanged.
\makeatletter\def\verbatim@font{\fontsize{8}{9.5}\selectfont\ttfamily}\makeatother

\title{Walkthrough Recording Runbook\\\large preflight, shots, verification, postflight}
\author{CS 454/554 Midterm Mini-Project --- companion to \code{docs/video-shotlist.md}}
\date{October 2026}

\begin{document}
\maketitle

\section*{How to use this}

Four minutes, six shots, one troubleshooting finding. Run the preflight in full; record only when every check passes. Each shot gives the block to paste, what the screen must show, the narration, and what to confirm before the next shot. The postflight restores the baseline and says how to reset for another take. Commands marked \emph{VM} run in an SSH session inside \code{cs554-vm}; commands marked \emph{host} run in Windows PowerShell. Counters and Pod names are live values: read what the screen shows, never the numbers printed here.

\paragraph{Windows to have open.} The public README in a browser (shots 1 and 6); \emph{Terminal 1}, an SSH session at \code{\textasciitilde/CS-554-MidTerm/compose}; \emph{Terminal 2}, an SSH session at \code{\textasciitilde/CS-554-MidTerm}; one host PowerShell window. Readable font at 1080p. Nothing personal visible: no taskbar, notifications, other tabs, or host paths.

\paragraph{Recorder.} Snipping Tool: \code{Win+Shift+S}, choose the video-camera icon, select the region that holds the terminals and the browser, enable the microphone, Start. Keep every window inside that region. Fallback: PowerPoint Record Screen, then a Teams screen-share recording. Record a five-second test clip first and play it back for audio.

\section{Preflight}

Run these in order; every one must pass. A failure is fixed with the postflight's restore steps, then the preflight is run again from the top. Do not record on a partial pass.

\paragraph{P1 --- VM running and addressed (host).}
\begin{verbatim}
$vmrun = "${env:ProgramFiles(x86)}\VMware\VMware Workstation\vmrun.exe"
& $vmrun -T ws list
& $vmrun -T ws getGuestIPAddress C:\VMs\cs554-vm\cs554-vm.vmx
\end{verbatim}
\expect{the list names \code{C:\textbackslash VMs\textbackslash cs554-vm\textbackslash cs554-vm.vmx}; the address is \code{192.168.197.x}. Write it down: shot 2 uses it.}

\paragraph{P2 --- Checkout at current \code{main} (VM).}
\begin{verbatim}
cd ~/CS-554-MidTerm
git status --porcelain
git fetch -q origin && git rev-parse --short HEAD origin/main
\end{verbatim}
\expect{no output from \code{status}; the two commit hashes are identical. If they differ: \code{git pull --ff-only origin main} and rerun.}

\paragraph{P3 --- Compose healthy (VM).}
\begin{verbatim}
docker inspect -f '{{.State.Health.Status}}' compose-web-1
curl -s localhost:8080/health
\end{verbatim}
\expect{\code{healthy}, then \code{\{"status":"ok"\}}.}

\paragraph{P4 --- Cluster at baseline (VM).}
\begin{verbatim}
kubectl get deploy cs554-midterm
kubectl get pods -l app=cs554-midterm
kubectl get pod client
kubectl get svc cs554-midterm -o jsonpath='{.spec.selector.app}'; echo
kubectl get endpointslices -l kubernetes.io/service-name=cs554-midterm -o json \
  | jq '[.items[].endpoints[]? | select(.conditions.ready == true)] | length'
for p in $(kubectl get pods -l app=cs554-midterm -o name); do
    kubectl exec "$p" -- test -f /tmp/ready && echo "$p readiness file present"
done
\end{verbatim}
\expect{Deployment \code{2/2}; two application Pods \code{1/1 Running} with \code{RESTARTS 0}; \code{client} \code{Running}; selector \code{cs554-midterm} (not \code{-typo}); ready endpoints \code{2}; both Pods report the readiness file.}

\paragraph{P5 --- No leftover forwarder (VM).}
\begin{verbatim}
pgrep -af 'kubectl.*port[-]forward' || echo "no forwarder"
ss -ltn '( sport = :18080 )'
\end{verbatim}
\expect{\code{no forwarder}; the \code{ss} table has a header line and nothing else.}

\paragraph{P6 --- Host reaches the service (host).}
\begin{verbatim}
curl.exe -i http://<address from P1>:8080/health
\end{verbatim}
\expect{\code{HTTP/1.1 200 OK} and \code{\{"status":"ok"\}}.}

\paragraph{P7 --- Recording setup.} Windows arranged inside the capture region, test clip recorded and played back with audible audio, this runbook open on another device.

\noindent\textbf{Go / no-go:} seven passes, start recording. Anything else: fix, rerun P1--P7, then record.

\section{Shots}

\subsection*{Shot 1 --- Open (0:00--0:20)}
\textbf{Screen:} the README in the browser: repository layout table and environment summary. No command.
\say{``This is one small service packaged as one image. Compose and Kubernetes both run inside this Linux VM on Windows, so the host, VM, container and application are separate layers.''}

\subsection*{Shot 2 --- Compose (0:20--1:10)}
\textbf{Screen:} Terminal 1, then the host PowerShell window.

\noindent Paste (VM, Terminal 1):
\begin{verbatim}
cd ~/CS-554-MidTerm/compose
docker compose ps
curl -s localhost:8080/
curl -s localhost:8080/
\end{verbatim}
\noindent Paste (host, with the address from P1):
\begin{verbatim}
curl.exe -s http://<address>:8080/
\end{verbatim}
\expect{\code{compose-web-1} \code{Up \ldots\ (healthy)} publishing \code{0.0.0.0:8080->8080/tcp}; three JSON responses with the same \code{hostname} and a \code{requestCount} that rises by one each time, the third one arriving from the host. Consecutive responses print on one line: the JSON has no trailing newline.}
\say{``The Windows host reaches the VM's NAT address directly; no port forwarding is needed on VMware. The hostname stays the same and the count rises by one with every request, because this counter belongs to that one application process.''}
\verify{the host response's count is exactly one more than the VM's last. If the host \code{curl.exe} fails, the address changed: confirm it with P1 and repeat the line.}

\subsection*{Shot 3 --- Kubernetes (1:10--2:00)}
\textbf{Screen:} Terminal 2.
\begin{verbatim}
cd ~/CS-554-MidTerm
kubectl get deploy,pods,svc
kubectl exec client -- sh -c 'for i in 1 2 3 4; do curl -s http://cs554-midterm/; echo; done'
\end{verbatim}
\expect{Deployment \code{2/2}; \code{client} and two application Pods \code{1/1 Running}; Service \code{cs554-midterm} of type \code{ClusterIP} on port \code{80}; four JSON lines whose \code{hostname} values are Pod names. Usually both Pods appear; a short sample that shows one hostname is not a failure (say so if it happens: ``a short random sample'').}
\say{``Kubernetes uses the same image, imported with \code{k3d image import}, and maintains two ready replicas. These in-cluster requests use Service DNS and ClusterIP routing; this sample reached both hostnames with independent counters, not guaranteed round-robin ordering.''}
\verify{all four requests returned JSON (no \code{curl:} error line).}

\subsection*{Shot 4 --- Pod replacement (2:00--2:40)}
\textbf{Screen:} Terminal 2. The first block selects a live Pod, deletes only that one, and watches; press \code{Ctrl+C} once the new Pod shows \code{1/1}, then paste the second block.
\begin{verbatim}
pod_to_delete=$(kubectl get pods -l app=cs554-midterm -o jsonpath='{.items[0].metadata.name}')
printf 'Selected Pod: %s\n' "$pod_to_delete"
kubectl delete pod "$pod_to_delete" --wait=false
kubectl get pods -w
\end{verbatim}
\begin{verbatim}
kubectl rollout status deploy/cs554-midterm --timeout=120s
kubectl get pods
\end{verbatim}
\expect{the selected Pod goes \code{Terminating}; a new Pod with a different generated name appears as \code{Pending} or \code{ContainerCreating}, then \code{Running 0/1}, then \code{1/1} within about ten seconds; \code{rollout status} prints \code{successfully rolled out}; the final list shows two application Pods \code{1/1 Running}, one of them seconds old.}
\say{``The ReplicaSet controller sees fewer Pods than the Deployment's desired two and creates a new named Pod automatically. Compose's manual restart reused its container and restarted the process; here Kubernetes replaces the Pod, and either new process loses its temporary counter.''}
\verify{two application Pods \code{1/1}, zero restarts, before starting shot 5. If the replacement has not reached \code{1/1} after 60 seconds, stop and diagnose with \code{kubectl describe pod <new-pod>}; do not continue.}

\subsection*{Shot 5 --- Troubleshooting finding (2:40--3:40)}
\textbf{Screen:} Terminal 2, no port-forward running. Paste the failure group, let the output settle, then paste the restore group, which is the README's ``Restore and verify'' block unchanged.

\noindent Failure group:
\begin{verbatim}
kubectl apply -f k8s/service-broken.yaml
kubectl get svc cs554-midterm -o jsonpath='{.spec.selector.app}'; echo
for attempt in $(seq 1 20); do
    count=$(kubectl get endpointslices -l kubernetes.io/service-name=cs554-midterm -o json \
      | jq '[.items[].endpoints[]? | select(.conditions.ready == true)] | length')
    if test "$count" -eq 0; then break; fi
    sleep 1
done
test "$count" -eq 0
failure_exit=0
kubectl exec client -- curl -sS --max-time 5 http://cs554-midterm/ || failure_exit=$?
printf 'Expected failed request exit: %s\n' "$failure_exit"
test "$failure_exit" -eq 7
kubectl get pods
kubectl get endpointslices -l kubernetes.io/service-name=cs554-midterm
\end{verbatim}
\expect{\code{service/cs554-midterm configured}; selector \code{cs554-midterm-typo}; \code{curl: (7) Failed to connect \ldots\ Could not connect to server} and \code{Expected failed request exit: 7}; both application Pods still \code{1/1 Running} with zero restarts; the EndpointSlice row shows \code{<unset>} under both \code{PORTS} and \code{ENDPOINTS}.}

\noindent Restore group:
\begin{verbatim}
kubectl apply -f k8s/service.yaml
for attempt in $(seq 1 20); do
    count=$(kubectl get endpointslices -l kubernetes.io/service-name=cs554-midterm -o json \
      | jq '[.items[].endpoints[]? | select(.conditions.ready == true)] | length')
    if test "$count" -eq 2; then break; fi
    sleep 1
done
test "$count" -eq 2 || { echo 'Two ready endpoints did not return'; exit 1; }
kubectl get svc cs554-midterm -o jsonpath='{.spec.selector}'
kubectl get endpointslices -l kubernetes.io/service-name=cs554-midterm -o wide
kubectl describe svc cs554-midterm | grep -E 'Selector|Endpoints'
recovered=
for attempt in $(seq 1 20); do
    if kubectl exec client -- curl -sS --max-time 5 http://cs554-midterm/; then recovered=$attempt; break; fi
    sleep 1
done
echo "recovered on attempt ${recovered:-none}"
kubectl get pods -l app=cs554-midterm --show-labels
\end{verbatim}
\expect{\code{service/cs554-midterm configured}; selector \code{\{"app":"cs554-midterm"\}}; the EndpointSlice lists two addresses; \code{Endpoints:} shows two \code{:8080} entries; possibly one or two \code{curl: (7)} lines, then a JSON response and \code{recovered on attempt N}; both Pods labeled \code{app=cs554-midterm}.}
\say{``Healthy Pods alone do not prove the Service works: its selector is a label query, and the typo leaves it with no matching endpoints. After correction, the endpoints can return before the node's forwarding rules converge; the recovery request is retried until it succeeds, here on attempt N.''}
\verify{the JSON response and the \code{recovered on attempt N} line are on screen with a number, not \code{none}. If it prints \code{none}: stop the recording, keep the output, and run the postflight restore; do not claim recovery on camera.}

\subsection*{Shot 6 --- Close (3:40--4:00)}
\textbf{Screen:} the README's ``Report and evidence'' links and the report instructions. No command.
\say{``The deployment and experiments are captured in the evidence files, and the technical report explains the comparison and OpenStack plan. The separate CS 554 readiness experiment is in \code{evidence/04-readiness-554.md} and section 6 of the report; it is not part of this on-camera selector demonstration.''}

\section{Postflight}

\subsection*{Restore (run after every take, finished or abandoned)}
\begin{verbatim}
cd ~/CS-554-MidTerm
kubectl apply -f k8s/service.yaml
kubectl rollout status deploy/cs554-midterm --timeout=120s
pkill -f 'kubectl.*port[-]forward' || true
\end{verbatim}
\expect{\code{service/cs554-midterm unchanged} (or \code{configured} if a take was abandoned mid-shot-5), \code{successfully rolled out}, and no error from \code{pkill}.}

\subsection*{Verify the baseline}
Rerun preflight checks P3, P4, and P5. All three must pass exactly as before: healthy Compose, \code{2/2}, two Pods \code{1/1} with zero restarts, selector \code{cs554-midterm}, two ready endpoints, both readiness files, no forwarder. The Pod names differ from before shot 4; that is expected and harmless.

\subsection*{Reset for another take (optional)}
Nothing needs resetting: counters simply continue, and the narration reads what the screen shows. If fresh counters are wanted anyway, this restarts the processes without changing any manifest or the Deployment:
\begin{verbatim}
cd ~/CS-554-MidTerm/compose && docker compose restart && cd ..
docker compose -f compose/compose.yaml ps
kubectl delete pod -l app=cs554-midterm
kubectl rollout status deploy/cs554-midterm --timeout=120s
\end{verbatim}
\expect{Compose back to \code{Up \ldots\ (healthy)} with the next request counting \code{1}; two new application Pods \code{1/1} with new names and counters at zero. Then run the full preflight again before recording.}

\subsection*{After the take}
Play the recording back for legible output and audible narration. Upload it unlisted, open the link in a private browser window to confirm it plays without sign-in, and send the URL to be added to the README and the Canvas note. Leave Compose, the cluster, and both replicas running until the submission is tagged.

\end{document}
